% Mise en page commune à tous les PDF du site
\usepackage{polyglossia}
\setdefaultlanguage{french}
\usepackage{fontspec}
\setmainfont{Latin Modern Roman}
\setsansfont{Latin Modern Sans}
\setmonofont{Latin Modern Mono}
\usepackage{amssymb}
\usepackage[most]{tcolorbox}
\usepackage{fancyhdr}
\usepackage{lastpage}
\usepackage{titlesec}
\usepackage{enumitem}
\setlist{itemsep=2pt, topsep=4pt}

% Couleurs (identiques au site)
\definecolor{cdef}{HTML}{3F51B5}
\definecolor{cthm}{HTML}{C62828}
\definecolor{cprop}{HTML}{00897B}
\definecolor{cdemo}{HTML}{6D4C41}
\definecolor{cdemoplus}{HTML}{7B1FA2}
\definecolor{cadmis}{HTML}{546E7A}
\definecolor{ctip}{HTML}{00838F}
\definecolor{cwarn}{HTML}{EF6C00}
\definecolor{cok}{HTML}{2E7D32}
\definecolor{cinfo}{HTML}{0277BD}
\definecolor{cetoile}{HTML}{F29100}
\definecolor{cprimaire}{HTML}{3949AB}

% Étoiles de difficulté (police qui contient le symbole ★)
\newfontfamily\policesymboles{DejaVu Sans}
\DeclareRobustCommand{\etoiles}[1]{\texorpdfstring{{\policesymboles\textcolor{cetoile}{#1}}}{#1}}

% Titres de sections
\titleformat{\section}{\Large\bfseries\color{cprimaire}}{\thesection.}{0.5em}{}
\titleformat{\subsection}{\large\bfseries}{}{0em}{}
\titleformat{\subsubsection}{\normalsize\bfseries}{}{0em}{}
\titlespacing*{\section}{0pt}{14pt}{6pt}
\titlespacing*{\subsection}{0pt}{12pt}{4pt}

% Encadrés : \begin{encadre}{couleur}{titre}
\newtcolorbox{encadre}[2]{
  enhanced, breakable,
  colback=#1!4!white, colframe=#1, colbacktitle=#1!12!white, coltitle=#1!70!black,
  boxrule=0.4pt, leftrule=3pt, arc=2pt, outer arc=2pt,
  fonttitle=\bfseries, title={#2},
  left=6pt, right=6pt, top=4pt, bottom=4pt,
  before skip=8pt, after skip=8pt
}


% En-tête et pied de page
\pagestyle{fancy}
\fancyhf{}
\fancyhead[L]{\small\sffamily\color{cprimaire}\textbf{Maths Pas à Pas}}
\fancyhead[R]{\small\sffamily\leftmark}
\fancyfoot[C]{\small\sffamily Page \thepage\ / \pageref*{LastPage}}
\renewcommand{\headrulewidth}{0.4pt}
\renewcommand{\sectionmark}[1]{}
\setlength{\headheight}{14pt}

% Pas d'indentation, paragraphes espacés
\setlength{\parindent}{0pt}
\setlength{\parskip}{5pt}

% Formules longues : autoriser la coupure
\allowdisplaybreaks

% Flèches et symboles Unicode tapés directement dans le texte
\usepackage{newunicodechar}
\newunicodechar{⇒}{\ensuremath{\Rightarrow}}
\newunicodechar{⇔}{\ensuremath{\Leftrightarrow}}
\newunicodechar{→}{\ensuremath{\rightarrow}}
\newunicodechar{←}{\ensuremath{\leftarrow}}
\newunicodechar{↔}{\ensuremath{\leftrightarrow}}
\newunicodechar{↗}{\ensuremath{\nearrow}}
\newunicodechar{↘}{\ensuremath{\searrow}}

% Tableaux : plus d'espace vertical (fractions dans les cellules)
\renewcommand{\arraystretch}{1.35}
\usepackage{array}
\setlength{\extrarowheight}{7pt}

% Figures : placées exactement à l'endroit prévu
\usepackage{float}
\floatplacement{figure}{H}
\newunicodechar{ᵉ}{\textsuperscript{e}}
\newunicodechar{✓}{\ensuremath{\checkmark}}
\newunicodechar{ʳ}{\textsuperscript{r}}
